\DocumentMetadata{
  pdfversion=2.0,
  pdfstandard=ua-2,
  lang=en-US,
  tagging=on,
  tagging-setup={math/setup=mathml-SE,tabsorder=structure}
}
\documentclass[11pt,letterpaper]{article}
\usepackage[margin=0.68in,includefoot]{geometry}
\usepackage{newcomputermodern}
\usepackage{amsmath,amscd,mathtools}
\providecommand{\square}{\mdlgwhtsquare}
\usepackage[hidelinks]{hyperref}
\hypersetup{
  pdftitle={Algebra PhD Exam, May 2007},
  pdfauthor={University of Florida Department of Mathematics},
  pdfsubject={Graduate mathematics examination},
  pdfdisplaydoctitle=true,
  bookmarksopen=true
}
\setcounter{secnumdepth}{0}
\setlength{\parindent}{0pt}
\setlength{\parskip}{0.28em}
\setlength{\emergencystretch}{3em}
\setlength{\leftmargini}{2.15em}
\setlength{\leftmarginii}{2.65em}
\setlength{\leftmarginiii}{3em}
\renewcommand{\labelenumi}{\textbf{\theenumi.}}
\pagestyle{empty}
\raggedbottom
\allowdisplaybreaks
\newenvironment{examproblems}[1]{%
  \begin{enumerate}%
  \setcounter{enumi}{#1}%
  \setlength{\topsep}{0.35em}%
  \setlength{\partopsep}{0pt}%
  \setlength{\itemsep}{0.48em}%
  \setlength{\parsep}{0.18em}%
}{\end{enumerate}}
\newenvironment{examsubparts}{%
  \begin{enumerate}%
  \setlength{\topsep}{0.18em}%
  \setlength{\partopsep}{0pt}%
  \setlength{\itemsep}{0.16em}%
  \setlength{\parsep}{0.1em}%
}{\end{enumerate}}
\newenvironment{examinstructions}{%
  \par\begingroup\itshape
}{%
  \par\endgroup\vspace{0.18em}
}
\makeatletter
\renewcommand\section{\@startsection{section}{1}{0pt}%
  {-1.25ex plus -.25ex minus -.1ex}%
  {0.55ex plus .1ex}%
  {\normalfont\large\bfseries}}
\renewcommand{\@maketitle}{%
  \newpage\null\vskip -1.2em%
  {\footnotesize\bfseries
    \href{https://www.ufl.edu/}{University of Florida}, \href{https://math.ufl.edu/}{Department of Mathematics}\par}%
  \vskip 0.18em%
  \tagstructbegin{tag=Title}%
    {\LARGE\bfseries\href{https://gma.math.ufl.edu/exams/algebra-phd/}{Algebra PhD Exam}, May 2007\par}%
  \tagstructend%
  \vskip 0.35em%
  \tagmcbegin{artifact}%
    \hrule height 1pt%
  \tagmcend%
  \vskip 0.55em%
}
\makeatother
\title{Algebra PhD Exam, May 2007}
\author{}
\date{}
\begin{document}
\maketitle
\thispagestyle{empty}
\begin{examinstructions}
Answer seven problems. You may quote results (within reason) as long as you state them clearly.
\end{examinstructions}
\begin{examproblems}{0}
\item
Give the definition of a normal algebraic field extension. Let~\(E\) be a normal algebraic extension of the field~\(F\) and \(g(x)\in F[X]\) an irreducible polynomial. Prove that the irreducible factors of \(g(x)\) in \(E[X]\) all have the same degree.
\item
Determine the Galois group of the polynomial \(x^5-2\in \mathbb{Q}[X]\).
\item
Define the terms natural transformation and natural isomorphism between functors~\(F\) and~\(G\) from a category~\(C\) to a category~\(D\). Show, giving full details, that \(M\mapsto R\otimes_R M\) defines a functor from the category of left~\(R\)-modules to itself, which is naturally isomorphic with the identity functor.
\item
Let~\(R\) be a ring with 1 and 
\[
0\longrightarrow A\longrightarrow B\longrightarrow C\longrightarrow 0
\]
 be an exact sequence of left unital~\(R\)-modules. Show that for any left~\(R\)-module~\(M\), the sequence of induced maps 
\[
0\longrightarrow \operatorname{Hom}_R(C,M)\longrightarrow \operatorname{Hom}_R(B,M)\longrightarrow \operatorname{Hom}_R(A,M)
\]
 is an exact sequence of abelian groups.
\item
Prove that the following conditions on a ring~\(R\) with 1 are equivalent:
\begin{examsubparts}
\item[\textnormal{(a)}]
Every unital left~\(R\)-module is projective.
\item[\textnormal{(b)}]
Every short exact sequence of left~\(R\)-modules splits.
\item[\textnormal{(c)}]
Every unital left~\(R\)-module is injective.
\end{examsubparts}
\item
Let~\(R\) be a commutative ring with 1 and~\(I\) the intersection of all prime ideals of~\(R\). Prove that for each \(a\in I\) there exists a natural number~\(n\) such that \(a^n=0\).
\item
Prove that the ring~\(R\) of \(n\times n\) matrices over a division ring is simple. Exhibit a minimal left ideal of~\(R\).
\item
Define a primitive ring and state and prove the Jacobson Density Theorem.
\item
Let~\(R\) be a Noetherian commutative ring with 1. Prove that \(R[X]\) is Noetherian.
\item
Give examples of the following. Explain your examples, stating clearly any results you use, but you do not need to prove all details.
\begin{examsubparts}
\item[\textnormal{(a)}]
A Noetherian ring which is not Artinian.
\item[\textnormal{(b)}]
A projective module which is not free.
\item[\textnormal{(c)}]
A Dedekind domain which is not a principal ideal domain.
\end{examsubparts}
\end{examproblems}
\end{document}
